\chapterhead{59}{ORR-20}{Capital Planning at Large Bank Holding Companies}{1}{2}

\lohead{59}{a}{Describe the Federal Reserve's Capital Plan Rule and explain the
seven principles of an effective capital adequacy process for bank holding
companies (BHCs) subject to the Capital Plan Rule.}

\orsub{The Federal Reserve's Capital Plan Rule}

The Federal Reserve ensures BHC stability through the following.

\begin{lettered}
\item \textbf{Capital Plan Rule:} requires BHCs with over \$50 billion in assets
to have a capital plan.
\item \textbf{CCAR (Comprehensive Capital Analysis and Review):} evaluates BHCs'
capital plans.
\end{lettered}

\orsub{The seven principles of an effective Capital Adequacy Process (CAP)}

\noindent\begin{minipage}{\textwidth}
\renewcommand{\arraystretch}{1.35}
\begin{tabularx}{\textwidth}{@{}L{0.26\textwidth}Y@{}}
\rowcolor{navy} \thd{Principle} & \thd{What it requires}\\
{\tabfont \textbf{a) Risk management}} & {\tabfont Comprehensive risk identification and control.}\\
\rowcolor{paleblue}
{\tabfont \textbf{b) Resource estimation}} & {\tabfont Clear definition and assessment of capital in stress scenarios.}\\
{\tabfont \textbf{c) Loss estimation}} & {\tabfont Process for estimating potential company-wide losses in stress conditions.}\\
\rowcolor{paleblue}
{\tabfont \textbf{d) Capital adequacy impact}} & {\tabfont Evaluation of loss impact on meeting capital requirements.}\\
{\tabfont \textbf{e) Capital policy}} & {\tabfont Well-defined strategy for capital goals, levels, distributions, and contingency plans.}\\
\rowcolor{paleblue}
{\tabfont \textbf{f) Internal controls}} & {\tabfont Robust validation, documentation, and audit mechanisms.}\\
{\tabfont \textbf{g) Oversight}} & {\tabfont Proactive governance by the board and senior management over the entire capital adequacy process.}\\
\end{tabularx}
\end{minipage}
\orfigcap{The seven CAP principles.}

% ---------------------------------------------------------------------
\lohead{59}{b}{Describe practices that can result in a strong and effective capital
adequacy process for a BHC in the following areas: risk identification; internal
controls, including model review and valuation; corporate governance; capital
policy, including setting of goals and targets and contingency planning; stress
testing and stress scenario design; estimating losses, revenues, and expenses,
including quantitative and qualitative methodologies; and assessing the impact of
capital adequacy, including risk-weighted asset (RWA) and balance sheet
projections.}

\orsub{Capital adequacy process}

\noindent\begin{minipage}{\textwidth}
\renewcommand{\arraystretch}{1.3}
\begin{tabularx}{\textwidth}{@{}YY@{}}
\rowcolor{navy} \thd{1) Risk identification} & \thd{2) Internal controls}\\
{\tabfont
a)~Identify all risk exposures across the BHC, including stress conditions,
changing economic environments, and on- and off-balance sheet items.\newline
b)~Regularly update the risk identification plan.\newline
c)~Capture ``other risks'' like compliance, reputational, and strategic risks.} &
{\tabfont
a)~Conduct thorough internal audits of capital planning data.\newline
b)~Maintain efficient Management Information Systems (MIS) for data collection and
analysis.\newline
c)~Implement a detailed documentation system for all CAP aspects.\newline
d)~Regularly review and validate all capital planning models independently.\newline
e)~Maintain a list of model inputs, assumptions, and adjustments.}\\
\rowcolor{navy} \thd{3) Corporate governance} & \thd{4) Capital policy}\\
{\tabfont
a)~Have a board with sufficient expertise to understand capital planning
processes.\newline
b)~Provide the board with comprehensive information on risks, losses, models,
etc.\newline
c)~Inform the board about stress scenarios and corrective actions.\newline
d)~Management should present capital plans to the board for approval.\newline
e)~Maintain detailed minutes of board meetings.} &
{\tabfont
a)~Clearly define capital goals, issuance, usage, and distributions.\newline
b)~Establish capital targets exceeding capital adequacy goals under stress.\newline
c)~Consider factors like future growth plans and economic conditions for
distributions.\newline
d)~Develop strong contingency plans for dealing with capital deterioration.}\\
\rowcolor{navy} \thd{5) Stress testing and scenario design} & \thd{6) Estimating losses, revenues, and expenses}\\
{\tabfont
a)~Design stress scenarios specific to the BHC's unique situation.\newline
b)~Use both internal models and expert judgment.\newline
c)~Stress testing models should consider all BHC risk exposures.\newline
d)~Scenarios should clearly explain how they address specific BHC weaknesses.} &
{\tabfont
a)~Use a mix of quantitative and qualitative methods for estimations.\newline
b)~Segment business lines and portfolios based on risk characteristics.\newline
c)~Employ sensitivity analysis to assess the impact of changing conditions.\newline
d)~Use conservative assumptions for estimations under normal and stress
conditions.}\\
\end{tabularx}
\end{minipage}
\orfigcap{Six of the seven CAP areas, with the practices expected in each.}

\orsubb{7) Loss estimation methods}

\begin{lettered}
\item Loss estimation methods should be theoretically sound and empirically valid.
\item Aggregate losses across business lines for firm-wide analysis using uniform
methods.
\item Use data that reflects the BHC's unique risk exposures for external data
sources.
\item Choose either economic loss (expected losses) or an accounting-based loss
approach.
\end{lettered}

\orsub{Assessing the impact of capital adequacy}

\orsubb{1) Operational risk}
\begin{lettered}
\item BHCs assess operational risk by examining the correlation with macroeconomic
factors, utilising a mix of techniques like regression models, modified loss
distribution approaches (LDA), and scenario analysis for robust loss estimates
under stress conditions.
\item When specific models or data are lacking, BHCs turn to scenario analysis to
cover a broader range of risks, justifying their scenario choices and performing
sensitivity analysis around selected confidence intervals for VaR estimates.
\end{lettered}

\orsubb{2) Market risk and counterparty credit risk}
\begin{lettered}
\item BHCs face market and counterparty credit risks in trading activities,
influenced by macroeconomic changes affecting the value of risk exposure and
counterparty creditworthiness.
\item They adopt probabilistic approaches to model potential portfolio losses,
providing evidence that these can capture more severe risk scenarios than
historical data. They also use deterministic approaches, applying a wide range of
stress scenarios to cover key exposures and explaining underlying assumptions and
mitigation measures.
\item Explicit consideration of counterparty default scenarios is made to address
significant risk exposures, with conservative assumptions about future risk
mitigation to account for limited actions in stress scenarios.
\end{lettered}

\orsubb{3) PPNR projection methodologies (Pre-Provision Net Revenue, less loss provisions)}
\begin{lettered}
\item Account for both present and future business and economic conditions in
stress scenario planning.
\item Ensure clear and consistent relationships between revenue, expenses, and
balance sheet projections. Maintain uniform assumptions across different
projections.
\item Base revenue, expense, and loss estimates on sound, well-documented
assumptions. Centralise projection oversight.
\item Utilise external data to supplement limited internal data.
\item Adapt net interest income projections to reflect balance sheet changes and
adjust for non-par value assets.
\item Clearly articulate projections based on product characteristics and
assumptions, like changes in deposit mix.
\item Explicitly link loss projections with net interest income, using models that
reflect loan portfolio behaviour.
\item Tailor non-interest income projections to the BHC's strategy and specific
business line risk exposures.
\item Show how off-balance sheet activities impact revenue projections.
\end{lettered}

\orsubb{4) Generating projections}

BHCs should have a well-defined process in place to develop projections of
revenues, expenses, losses, on- and off-balance sheet items, and risk-weighted
assets in an enterprise-wide scenario analysis. Projections should be based on
sound underlying assumptions, interactions, and factors (main drivers of change),
and the estimates should be scrutinised, documented, and reported.
